\chapter{Experimental procedures}
\label{app:procedures}

\section{Chemical cleaning procedure}
\label{sec:cleaning-procedure}
The same procedure was used for the commercial (Campaign~1) and extensive (Campaign~2) cleans. Chemical quantities, solution volumes, measured pH, soak times and agitation for each protocol are given in Table~\ref{tab:cleaning-protocols}. All solutions were prepared with Milli-Q water, all soaks were carried out at room temperature (approximately \qty{22}{\degreeCelsius}) without heating, and pH was measured with a calibrated pH meter immediately after each solution was prepared. Each membrane sheet was cleaned separately.

\subsection*{Alkaline (caustic) clean with chelant}
\begin{enumerate}[nosep]
  \item The membrane sheet was cut to size (approximately \qty{300}{\square\centi\metre}; \qty{30}{\centi\metre} $\times$ \qty{10}{\centi\metre} for the extensive clean).
  \item The required volume of Milli-Q water was measured out. NaOH pellets were weighed and dissolved in a portion of this water with manual stirring until fully dissolved.
  \item \ce{Na2EDTA} was weighed and dissolved into the same solution.
  \item The pH of the combined solution was measured and recorded.
  \item The membrane sheet was fully submerged in the alkaline solution, feed side up, and soaked for the specified time. In the commercial clean, the sheet was laid flat in a clear dish and soaked statically. In the extensive clean, the vessel was covered with aluminium foil and the solution was stirred with a magnetic stirrer at \qty{100}{rpm} throughout the soak.
\end{enumerate}

\subsection*{Intermediate rinse}
\begin{enumerate}[resume, nosep]
  \item The membrane was rinsed continuously with RO water for \qty{10}{\minute} to remove residual alkaline solution, which would otherwise neutralise the acid solution on contact.
\end{enumerate}

\subsection*{Acid clean}
\begin{enumerate}[resume, nosep]
  \item \qty{32}{\percent} w/w HCl was added to the required volume of Milli-Q water and mixed.
  \item The pH of the solution was measured and recorded.
  \item The membrane sheet was fully submerged in the acid solution and soaked for the specified time, statically in the commercial clean and covered and stirred at \qty{100}{rpm} in the extensive clean.
\end{enumerate}

\subsection*{Final rinse and storage}
\begin{enumerate}[resume, nosep]
  \item The membrane was rinsed continuously with RO water until free of cleaning agents.
  \item Cleaned sheets were either cut directly into rig coupons or used for IP (Section~\ref{sec:ip-method}). Sheets and coupons were stored in Milli-Q water in transparent airtight bags in a refrigerator until testing or analysis.
\end{enumerate}